\section{Approach}
\label{sec:approach}

Figure~\ref{fig:overview} presents the workflow of \tool.
Our core concept is to build an adaptive, structure-aware fuzzing base for binary asset formats (\fastgltf).
We then systematically extend this architecture to interactive multimedia systems (\raylib) and asynchronous web servers (\drogon) using domain-specialized headless in-memory shims.
We do not claim standard compiler instrumentation, libFuzzer's in-process loop, or AddressSanitizer runtime checks as new.
Our design contribution lies in the cross-domain fuzzer adaptation, the decoupled bandit mutation scheduler, and the zero-socket in-memory network mocking layer.

\begin{figure}[t]
  \centering
  %% Overview of the unified fuzzer architecture and cross-domain extensions.
%% Grayscale-safe, clean modern styling, publication-grade alignment with matched heights.
\begin{tikzpicture}[
    >=Stealth, font=\scriptsize,
    titleheader/.style={
      font=\tiny\bfseries, align=center, text width=32.5mm,
      fill=black!10, draw=black!60, rounded corners=2pt, inner sep=2pt
    },
    stagecard/.style={
      draw=black!70, rounded corners=2pt, align=left, inner sep=3pt,
      text width=32.5mm, minimum height=16.5mm, fill=white, font=\tiny, line width=0.5pt
    },
    oraclecard/.style={
      draw=black!75, rounded corners=2pt, align=center, inner sep=3.5pt,
      text width=113mm, minimum height=9mm, fill=black!7, font=\tiny, line width=0.6pt
    },
    flow/.style={->, line width=0.6pt, draw=black!85},
    back/.style={->, line width=0.6pt, dashed, draw=black!75},
    extlabel/.style={
      font=\tiny\bfseries, fill=white, inner sep=1pt, rounded corners=1pt,
      draw=black!35
    }
  ]
  % Stage 1: Base Engine (fastgltf)
  \node[titleheader] (h1) at (0, 0) {(1) Base Engine (\fastgltf)};
  \node[stagecard, below=1.8mm of h1] (s1) {%
    $\triangleright$ Structured AST Mutators\\
    $\triangleright$ Binary Geometry Mutators\\
    $\triangleright$ Discounted-UCB1 Bandit%
  };

  % Stage 2: Multimedia Extension (raylib)
  \node[titleheader, right=7.5mm of h1] (h2) {(2) Multimedia (\raylib)};
  \node[stagecard, below=1.8mm of h2] (s2) {%
    $\triangleright$ 37 Subsystem Binaries\\
    $\triangleright$ Headless \texttt{rlsw} Rasterizer\\
    $\triangleright$ Virtual VFS \& Audio Pumps%
  };

  % Stage 3: Network Extension (drogon)
  \node[titleheader, right=7.5mm of h2] (h3) {(3) Network I/O (\drogon)};
  \node[stagecard, below=1.8mm of h3] (s3) {%
    $\triangleright$ 19 Stateful Target Binaries\\
    $\triangleright$ \texttt{MockTcpConnection} (0-sock)\\
    $\triangleright$ Tail Slicing \& Sandboxing%
  };

  % Combined Oracle & Feedback Layer
  \node[oraclecard, below=5.5mm of s2] (o) {%
    \textbf{Sanitizer Oracle (ASan / UBSan) \& Linker Coverage Interception (\texttt{-Wl,--wrap})}\\[1.5pt]
    Dual-channel edge telemetry $\bullet$ Crash triage engine $\bullet$ 16 discovered security defects%
  };

  % Horizontal extension flows between headers
  \draw[flow] (h1.east) -- (h2.west) node[midway, extlabel] {extend};
  \draw[flow] (h2.east) -- (h3.west) node[midway, extlabel] {extend};

  % Downward execution paths into the oracle
  \draw[flow] (s1.south) -- (s1.south |- o.north);
  \draw[flow] (s2.south) -- (o.north);
  \draw[flow] (s3.south) -- (s3.south |- o.north);

  % Feedback loop from coverage interceptor back to the bandit scheduler
  \draw[back] (o.west) -- ++(-4.5mm, 0) |- node[pos=0.25, left, font=\tiny\bfseries, align=right] {bandit\\reward} (s1.west);

  % Background card frames behind each stage
  \begin{scope}[on background layer]
    \node[draw=black!25, fill=black!1, rounded corners=3pt, line width=0.4pt,
          fit=(h1) (s1), inner sep=2.5pt] {};
    \node[draw=black!25, fill=black!1, rounded corners=3pt, line width=0.4pt,
          fit=(h2) (s2), inner sep=2.5pt] {};
    \node[draw=black!25, fill=black!1, rounded corners=3pt, line width=0.4pt,
          fit=(h3) (s3), inner sep=2.5pt] {};
  \end{scope}

\end{tikzpicture}

  \caption{Overview of \tool. The system starts with the base structure-aware engine for \fastgltf (\circled{1}), extends to multimedia subsystems in \raylib via headless rendering mocks (\circled{2}), and extends to network servers in \drogon via in-memory socket abstractions (\circled{3}). Coverage rewards update the bandit mutation scheduler.}
  \label{fig:overview}
\end{figure}

\subsection{Step 1: Base Engine and Adaptive Mutation (\fastgltf)}
\label{sec:base_engine}
We began by designing \tool for the \fastgltf 3D asset library.
Because glTF 2.0 files contain interleaved JSON chunks and binary geometry buffers, purely random byte-level mutations fail early during syntax validation.
We built a hybrid mutation engine combining four distinct operators:
(i)~\emph{structural field mutations} that modify JSON accessor indices, buffer views, and component counts;
(ii)~\emph{binary geometry mutators} that modify vertex coordinates, normals, and indices;
(iii)~\emph{integer boundary overrides} that inject extreme values such as \texttt{SIZE\_MAX}, 0, and negative strides; and
(iv)~\emph{standard bit and byte flips}.

To steer these operators adaptively, we implemented the UCB1 multi-armed bandit algorithm (Equation~\ref{eq:ucb1}).
Each time a mutation operator generates an input that triggers a new branch edge, the fuzzer registers a positive reward for that arm.
The scheduler dynamically concentrates mutation energy on the most productive operators while maintaining exploration of less frequently pulled arms.

\subsection{Step 2: Multimedia Extension and Subsystem Decomposition (\raylib)}
\label{sec:multimedia_ext}
Next, we extended the architecture to \raylib, a C multimedia and game development library.
Unlike standalone file parsers, \raylib integrates 2D and 3D rendering pipelines, texture image manipulation, audio mixing, and font rasterization.
Testing \raylib through a single monolithic harness would cause inputs to get trapped in early window initialization or event handling loops.

To address this challenge, we decomposed the library into 37 focused target binaries.
Each target isolates a specific subsystem: the Image VM (52 operators), Render VM (software rasterizer \texttt{rlsw}), audio streams, or event replay.
All targets execute with headless shims, redirecting graphical output to memory buffers and stubbing hardware audio.
Figure~\ref{fig:subsystems}(a) illustrates the headless multimedia architecture for \raylib.

\subsection{Step 3: Network Extension and In-Memory Execution (\drogon)}
\label{sec:network_ext}
Finally, we extended \tool to \drogon, a high-throughput C++ asynchronous web framework.
Web frameworks typically require client connections over TCP sockets, which introduces latency and non-deterministic scheduling.

We solved this through three design elements.
First, \texttt{MockTcpConnection} pushes buffers directly into parser callbacks in memory without socket allocations.
Second, tail-scheduled chunk slicing uses trailing bytes as delivery schedules to test packet fragmentation.
Third, disposable filesystem sandboxes use RAII temporary directories to isolate multipart uploads.
Figure~\ref{fig:subsystems}(b) details the in-memory architecture for \drogon.

\begin{figure}[t]
  \centering
  %% Detailed architecture diagrams for raylib and drogon fuzzers.
%% Side-by-side panels: (a) raylib headless architecture, (b) drogon in-memory architecture.
%% Symmetrical grid alignment, publication-grade layout with balanced heights and container cards.
\begin{tikzpicture}[
    >=Stealth, font=\scriptsize,
    titlebox/.style={
      font=\scriptsize\bfseries, align=center, text width=54mm,
      fill=black!10, draw=black!65, rounded corners=2pt, inner sep=2.5pt
    },
    tier1/.style={
      draw=black!70, rounded corners=2pt, align=center, inner sep=2pt,
      text width=53mm, minimum height=5.5mm, fill=white, font=\tiny, line width=0.5pt
    },
    tier2/.style={
      draw=black!75, rounded corners=2pt, align=left, inner sep=2.5pt,
      text width=53mm, minimum height=18.5mm, fill=white, font=\tiny, line width=0.6pt
    },
    tier3/.style={
      draw=black!75, rounded corners=2pt, align=center, inner sep=2.5pt,
      text width=53mm, minimum height=8mm, fill=black!5, font=\tiny, line width=0.6pt
    },
    tier4/.style={
      draw=black!80, rounded corners=2pt, align=center, inner sep=2pt,
      text width=53mm, minimum height=5.5mm, fill=black!12, font=\tiny, line width=0.6pt
    },
    arr/.style={->, line width=0.6pt, draw=black!80}
  ]

  % ================= PANEL A: RAYLIB =================
  \node[titlebox] (r_title) at (0, 0) {(a) \raylib Headless Multimedia Fuzzer};

  \node[tier1, below=1.5mm of r_title] (r_in) {%
    \textbf{Raw Byte Stream} $\to$ \texttt{fz\_rd} Parameter Reader%
  };

  \node[tier2, below=1.5mm of r_in] (r_shim) {%
    \textbf{Headless Runtime Shims (\texttt{PLATFORM\_MEMORY}):}\\[1.5pt]
    $\triangleright$ \texttt{rlsw} software rasterizer \& virtual framebuffer\\
    $\triangleright$ In-memory VFS via \texttt{SetLoadFileCallback}\\
    $\triangleright$ Capped allocator (\texttt{RL\_MALLOC $\to$ fz\_malloc})\\
    $\triangleright$ Synchronous audio mixer pump (no sound hardware)%
  };

  \node[tier3, below=1.5mm of r_shim] (r_tgt) {%
    \textbf{37 Isolated Subsystem Target Binaries}\\[1pt]
    Image VM (52 ops) $\bullet$ Render VM $\bullet$ Audio $\bullet$ Fonts $\bullet$ Models%
  };

  \node[tier4, below=1.5mm of r_tgt] (r_ora) {%
    \textbf{Sanitizer Oracle (ASan / UBSan) \& Crash Triage}%
  };

  \draw[arr] (r_in) -- (r_shim);
  \draw[arr] (r_shim) -- (r_tgt);
  \draw[arr] (r_tgt) -- (r_ora);

  % ================= PANEL B: DROGON =================
  \node[titlebox] (d_title) at (60.5mm, 0) {(b) \drogon In-Memory Network Fuzzer};

  \node[tier1, below=1.5mm of d_title] (d_in) {%
    \textbf{Raw Stream} $\to$ Tail Slicer (trailing 4B schedule)%
  };

  \node[tier2, below=1.5mm of d_in] (d_shim) {%
    \textbf{In-Memory Network Shims (Zero OS Sockets):}\\[1.5pt]
    $\triangleright$ \texttt{MockTcpConnection} captures virtual interface\\
    $\triangleright$ Thread-local \texttt{EventLoop} binding (race-free execution)\\
    $\triangleright$ Edge-biased slicing into \texttt{trantor::MsgBuffer}\\
    $\triangleright$ RAII \texttt{IterationSandbox} directory containment%
  };

  \node[tier3, below=1.5mm of d_shim] (d_tgt) {%
    \textbf{19 Stateful Network Target Binaries}\\[1pt]
    HttpRequest $\bullet$ HttpResponse $\bullet$ Multipart $\bullet$ WebSocket%
  };

  \node[tier4, below=1.5mm of d_tgt] (d_ora) {%
    \textbf{Sanitizer Oracle (ASan / UBSan) \& Crash Triage}%
  };

  \draw[arr] (d_in) -- (d_shim);
  \draw[arr] (d_shim) -- (d_tgt);
  \draw[arr] (d_tgt) -- (d_ora);

  % Background card frames
  \begin{scope}[on background layer]
    \node[draw=black!25, fill=black!1, rounded corners=3pt, line width=0.4pt,
          fit=(r_title) (r_in) (r_shim) (r_tgt) (r_ora), inner sep=3.5pt] {};
    \node[draw=black!25, fill=black!1, rounded corners=3pt, line width=0.4pt,
          fit=(d_title) (d_in) (d_shim) (d_tgt) (d_ora), inner sep=3.5pt] {};
  \end{scope}

\end{tikzpicture}

  \caption{Subsystem fuzzer architectures: (a)~\raylib headless multimedia fuzzer executing over virtual memory framebuffers and audio streams; and (b)~\drogon in-memory web framework fuzzer driving asynchronous parsers via \texttt{MockTcpConnection} and tail-scheduled chunk slicing.}
  \label{fig:subsystems}
\end{figure}

\subsection{Step 4: Independent Coverage Interception via Linker Wrapping}
\label{sec:linker_wrapping}
To track bandit rewards independently of libFuzzer's internal state, the mutator requires edge discovery signals.
However, in modern toolchains such as Clang 18, compiling with \texttt{-fsanitize=fuzzer} provides a strong internal definition of \texttt{\_\_sanitizer\_cov\_trace\_pc\_guard}.
Defining our own hook caused linker duplicate-symbol collisions.

We resolved this conflict by passing GNU linker symbol wrapping flags (\texttt{-Wl,--wrap=\_\_sanitizer\_cov\_trace\_pc\_guard} and \texttt{-Wl,--wrap=\_\_sanitizer\_cov\_trace\_pc\_guard\_init}).
Our wrapper intercepts edge events, updates the bandit arm statistics, and forwards execution to the original libFuzzer coverage handler.
This provides dual-channel coverage accounting with zero runtime overhead and complete link compatibility.
We implemented \tool in 4{,}350 lines of C and C++17 compiled under Clang 18.1.3 on Linux x86\_64 (Ubuntu 24.04 LTS).
